% ================= CHAPTER 1 =================
\chapter{PENDAHULUAN}

\section{Latar Belakang}
Latar belakang mencerminkan alur pemikiran yang dinamis, mengapa suatu fenomena (gejala alam, gejala sosial) yang dijumpai dapat menggugah niat atau panggilan untuk melakukan percobaan penelitian \parencite{dewi2022}. Sehubungan dengan ini, peneliti harus merasa yakin, bahwa fenomena yang dijumpai itu merupakan masalah yang masih aktual dan relevan dengan keadaan sekarang [cite: 250]. Hal tersebut dapat diperoleh dengan jalan mencari informasi dari kepustakaan yang terbaru atau berkonsultasi dengan pakar ilmiah yang berkaitan dengan disiplin ilmu yang bersangkutan [cite: 251].

Dari pihak peneliti, pengungkapan bagian ini dapat didasarkan atas pertanyaan-pertanyaan berikut:
\begin{itemize}
    \item Apa yang telah diketahui, teoritis maupun faktual, dari masalah yang diteliti?
    \item Adakah permasalahan disitu, apakah ada keraguan yang terdapat pada permasalahan itu?
    \item Apa hal yang menarik dari masalah yang diteliti?
    \item Apakah mungkin secara teknis masalah itu diteliti?
\end{itemize}
Pada prinsipnya latar belakang percobaan penelitian dapat menjawab pertanyaan mengapa masalah itu diteliti.

\section{Perumusan Masalah}
Identifikasi masalah merupakan penjabaran dari masalah, menjadi beberapa sub masalah yang spesifik dan biasanya dirumuskan dalam kalimat tanya.

\textbf{Contoh:} \\
Berdasarkan uraian latar belakang yang telah dipaparkan, maka rumusan masalah yang diangkat pada penelitian ini adalah sebagai berikut:
\begin{enumerate}
    \item Bagaimana struktur sistem nanokomposit yang dapat menunjukkan efisiensi two photon absorbtion (TPA) yang baik agar proses ini dapat dihasilkan pada intensitas cahaya yang lebih rendah?
    \item Bagaimana mekanisme TPA dan karakteristik proses tersebut pada sistem nanohibrid?
    \item Bagaimana karakteristik TPA pada semiconductor quantum dot (SQD) yang dimodelkan lebih kompleks dari two level system (TLS)?
\end{enumerate}

\section{Tujuan Penelitian}
Tujuan penelitian adalah jawaban rumusan masalah. Dengan mengacu kepada contoh perumusan identifikasi masalah, maka tujuan penelitiannya adalah sebagai berikut:

\textbf{Contoh:} \\
Penelitian ini bertujuan untuk mempelajari proses TPA pada nanokomposit SQD-MNP. Secara khusus tujuan tersebut dapat dirinci sebagai berikut:
\begin{enumerate}
    \item Memperoleh konfigurasi struktur sistem nanokomposit dengan efisiensi TPA yang baik sehingga proses tersebut dapat dieksitasi pada intensitas cahaya yang lebih rendah.
    \item Mengungkap prinsip dan mekanisme TPA serta karakteristik proses tersebut pada sistem nanokomposit.
    \item Memahami karakteristik TPA pada SQD untuk sistem yang kompleks dari TLS.
\end{enumerate}

\section{Manfaat Penelitian}
Sub bab ini merupakan penajaman spesifikasi sumbangan hasil percobaan penelitian terhadap nilai manfaat praktis, juga sumbangan ilmiahnya bagi perkembangan ilmu.

\textbf{Contoh:} \\
Manfaat pada penelitian ini adalah diperolehnya informasi mengenai metode sintesis material fotokatalis PP-TiO$_2$ serta pengaruh material ini pada proses fotodegradasi senyawa organik air gambut. Walaupun penelitian ini masih dalam cakupan skala laboratorium, hasil dari riset ini berpotensi diimplementasikan pada sektor industri pengolahan air bersih.
